\subsection{Gap and Positioning of This Work}

Existing research has demonstrated significant progress in LLM-based multi-agent systems and their coordination mechanisms, including architectures for reasoning and collaborative problem solving.
However, the evaluation of such systems remains largely grounded in well-defined benchmark setting with explicit ground-truth signals and objective correctness criteria.
This creates a fundamental mismatch between the capabilities of modern LLM-MAS and the evaluation protocols used to assess them.

This limitation is particularly pronounced in urban planning, where problems are underspecified, admit multiple valid solutions and require trade-offs between competing objectives. 
Existing approaches either reduce problems to structured subtasks or rely on heuristic evaluation that do not capture the full complexity of generated solutions.

In this work, we address this gap by proposing:

\begin{itemize}
    \item \textbf{FP2MP}: an LLM-MAS architecture for solving ill-defined urban planning problems inspired by LbMAS architecture \cite{han2025exploring}.
    \item \textbf{FP2MP-Eval}: a multi-dimensional LLM-as-a-judge evaluation framework for assessing reasoning quality, problem structuring, and reflective capabilities of LLM-MAS in ill-defined settings.
\end{itemize}